\begingroup
\normalsize
\begin{flushleft}
\bfseries\abstractname
\end{flushleft}
\noindent
This study evaluates signal reconstruction through Gaussian noise and a
binary symmetric channel using standard signal-processing methods. DFT analysis selects
16 kHz sampling for the supplied two-second multitone, retaining all four
components while reducing samples by two thirds. The baseline combines a
385-tap Hann multiband FIR, 14-bit endpoint quantisation, Hamming(7,4) and
periodic Fourier reconstruction. Purpose-seeded channels use separate
random streams under a documented project mapping.

Independent four-configuration comparisons give the lowest anchor-point
mean NMSE, $\WPersonalBoth$, when filtering and coding are both enabled.
Filtering becomes harmful in Gaussian-only 35/40 dB checks. Seven of 92
R4 means miss the planned 5\% relative interval half-width target (maximum
\WMaxRelativePercent\%); some high-SNR contrasts remain unresolved.

Changing to 16 bits improves the declared main-domain objective but raises
traffic from 392 to 448 kbit/s, with unresolved anchor-point benefit.
Under a separate 320 kbit/s cap, Hamming(15,11) meets the fixed anchor-point
mean-quality rule at 305.46 kbit/s, with reduced exact-payload reliability.
These conclusions are conditional on the tested signal, severity weights
and logical-rate model. Reproducible code, saved measurements and independent
numerical-method checks support the results; the numerical assumptions and limitations are stated explicitly.
\par
\endgroup
